\section{内部の仕組み}\label{sec:internals}
この章は，スタイルを改造・保守する人のための短い説明です。利用者は読まなくてかまいません。詳しくは，\file{jsr.bbx}，\file{jsr.cbx}，\file{jsr.dbx} のコメントと，リポジトリの \file{tests/README.md}，\file{docs/DECISIONS.md} にあります。

\subsection{原典を代表にする方式}
訳書・再録（\bibfield{related} ＋ \bibfield{relatedtype} が \texttt{translationof} または \texttt{reprintof}）は，単独の項目にせず，原典の項目に添えて出します。原典が引用されていないときは，biber が作る原典の複製（clone）が代表になります。

\begin{itemize}
\item \cmd{DeclareStyleSourcemap} で，related のオプションと並べ替えの役割（\bibfield{jsrsortrole}）を付けます。
\item \cmd{AtDataInput}（\cmd{jsr@datainput}）で，読み込んだ項目に添え物の印を付けます。
\item \cmd{blx@bbl@enddlist} の直前に \cmd{jsr@finishlist} を差し込み，添え物を印刷リストから除いて，extradate（a，b）を振り直します。
\item 添え物は，\texttt{jsr:work}，\cmd{jsr@printattach}，\texttt{jsr:attachbody} で出します。
\end{itemize}
\texttt{blx@dlist@entry@…}，\texttt{blx@data@…}，\cmd{blx@bbl@enddlist} は biblatex の内部名です。biblatex を更新して壊れたら，まずここを疑ってください。

\subsection{sortname の印}
\bibfield{sortname} は bbl に書き出されません（biber が並べ替えにだけ使います）。TeX 側で有無を調べたいときは，sourcemap で印のフィールド（\bibfield{jsrnosortname}，\bibfield{jsrsortrole}。\file{jsr.dbx} で宣言）を付けます。

\subsection{署名のない記事の除外}
\bibtype{article} のドライバを使う型（\bibtype{article}，\bibtype{misc}，\bibtype{software}，\bibtype{patent}）で \bibfield{author} がないものには，sourcemap で \bibfield{options} に \texttt{skipbib} を加えます（\bibfield{options} が既にあるときは追記します）。ドライバが何も出さないだけだと，文献リストに空の項目（空行）が残るためです。\texttt{skipbib} の項目も，ラベルは通常どおり付くので，文献注と extradate には影響しません。

新聞・雑誌（\bibfield{entrysubtype} または互換の \bibfield{type} が \texttt{newspaper}／\texttt{magazine}）以外は，警告用の印（\bibfield{jsrnoauthor}）も付け，\cmd{AtDataInput} の処理（\cmd{jsr@datainput}）が警告を出します。\texttt{skipbib} の項目は文献リストに組まれないため，\cmd{AtEveryBibitem} では警告が出せないからです。\bibfield{crossref}／\bibfield{xref} がある項目は，author を親から継承するかもしれないので対象外です。

値が空白だけのフィールド（\texttt{author = \{\}}）は，ほかの規則より前に，sourcemap で消します（\bibfield{author}，\bibfield{editor}，\bibfield{translator}，\bibfield{sortname}）。biber の \texttt{notfield} は「存在するか」を調べるので，空の値が残ると，署名のない記事の除外や sortname の警告が効かないためです。

\subsection{旧形式の読み替え}
\bibfield{webpagetitle} と \texttt{type=onlinereport} は，sourcemap で読み替えます。このとき，警告用の印（\bibfield{jsrlegacy}）を付けておき，組版時に警告を出します。

\subsection{言語の判定}
\bibfield{langid} が読み込まれているかは，プリアンブル終了時の言語一覧（\cmd{jsr@doclangs}）で判定します。babel の \cmd{bbl@loaded}，\cmd{bbl@ini@loaded}，polyglossia の \cmd{xpg@bloaded}，\cmd{xpg@loaded} を読むので，これらの内部名が変わったら壊れます。一覧にない言語の項目は，\bibfield{langid} を消して american で組みます。

\subsection{実装上の注意}
\begin{itemize}
\item \cmd{savebox} の中で \cmd{printnames} を使うとき（発行元の判定，カタカナ名の判定）は，\cmd{jsr@savepunct}／\cmd{jsr@restorepunct} で句読点の状態を退避します。しないと，保留中の区切りが消えます。
\item ループで \cmd{do} を使い回すと，入れ子で壊れます。新しいループには専用のハンドラを使います。
\item リストが空になりうる箇所では，etoolbox のループの代わりに \cmd{jsr@forlistcs}／\cmd{jsr@dolistcs} を使います。未定義のリストを渡すと，\cmd{relax} が項目として処理されるためです。
\end{itemize}
